\chapter{Des individus uniques}\label{ch:g9:unique-individuals}

Sept milliards de personnes, et pas deux pareilles --- ni
dans une famille, ni dans une cour de récréation, ni dans
toute l'histoire, les vrais jumeaux à moitié exceptés. Trois
chapitres ont maintenant assemblé toute la machinerie de ce
fait. Ce court chapitre de couronnement déroule l'argument
une fois, d'un bout à l'autre, et en encaisse les
conséquences~: pour la variété de l'\omterm{def:g5:biodiversity:species}{espèce}, pour notre
lecture des titres de journaux sur «~le génétique~», et pour
ce que l'unicité veut dire et ne veut pas dire.

\section{Le théorème de l'unicité}

\begin{proposition}[Pourquoi deux personnes ne se ressemblent jamais]\label{prop:g9:unique-individuals:theorem}
Assemble les trois mécanismes~:
\begin{enumerate}
  \item \emph{le brassage de la réduction de moitié}~: les
        \omterm{def:g8:sexual-reproduction:gametes}{gamètes} de chaque parent portent un partenaire par
        paire, distribué indépendamment --- $2^{23}$
        demi-jeux par parent
        (\cref{rem:g9:chromosomes:whyhalve}), et en vérité
        bien davantage, puisque les partenaires appariés
        \emph{échangent aussi des segments} pendant la
        division réductrice, coupant le jeu en dessous du
        niveau des \omterm{def:g9:chromosomes:chromosomes}{chromosomes} entiers~;
  \item \emph{la loterie de la \omterm{def:g5:flower-to-fruit:steps}{fécondation}}~: quelle
        \omterm{def:g8:reproductive-systems:male}{cellule mâle} parmi les millions
        (\cref{ex:g8:reproductive-systems:sperm}) rencontre
        la \omterm{def:g4:animal-reproduction:sexual}{cellule œuf} de quel mois multiplie entre eux les
        deux espaces~;
  \item \emph{le filet de la \omterm{prop:g9:genes-and-dna:explains}{mutation}}~: chaque génération
        ajoute quelques orthographes neuves qui lui sont
        propres (\cref{prop:g9:genes-and-dna:explains}).
\end{enumerate}
L'espace des combinaisons écrase le nombre d'êtres humains
qui ont jamais vécu~: chaque \omterm{def:g5:flower-to-fruit:steps}{fécondation} distribue un jeu
qui n'a jamais été distribué et ne le sera jamais plus.
\end{proposition}
\admitted

\begin{example}[Les jumeaux, et les limites de l'exception]\label{ex:g9:unique-individuals:twins}
Les vrais jumeaux partagent une seule distribution
(\cref{exo:g8:from-egg-to-birth:10}) --- et même eux
divergent~: des vies séparées écrivent des \omterm{def:g9:heredity-and-traits:traits}{caractères} acquis
différents (\cref{prop:g9:heredity-and-traits:sorting}), des
histoires de \omterm{def:g8:nervous-communication:synapse}{synapses} différentes
(\cref{ex:g8:nervous-communication:juggler}), et même
quelques \omterm{prop:g9:genes-and-dna:explains}{mutations} tardives différentes dans les \omterm{def:g5:first-look-at-cells:cell}{cellules} de
chaque corps. La seule fissure de l'unicité se referme
d'elle-même~: après la distribution, vivre différencie.
\end{example}

\begin{remark}[Uniques, et pour l'essentiel semblables]\label{rem:g9:unique-individuals:alike}
Tiens les deux vérités~: les textes d'\omterm{def:g9:genes-and-dna:dna}{ADN} de deux humains
quelconques s'accordent sur plus de $99$ pour cent --- la
construction partagée de l'\omterm{def:g5:biodiversity:species}{espèce} unique de
\cref{prop:g5:first-look-at-cells:principle} --- et les
désaccords, éparpillés dans trois milliards de lettres,
suffisent quand même à une unicité stricte. «~Tous
différents~» et «~tous parents~» sont le même fait à deux
grossissements (\cref{met:g9:genes-and-dna:levels})~; tout
usage jamais fait de la première vérité contre la seconde
était de la mauvaise \omterm{def:g1:living-or-not:living}{biologie} avant d'être de la mauvaise
morale.
\end{remark}

\section{L'unicité au travail}

\begin{example}[L'identification par l'ADN]\label{ex:g9:unique-individuals:forensics}
L'unicité est vérifiable, et les tribunaux la vérifient~:
un nombre suffisant de positions variables d'une personne,
lues à partir d'une trace de \omterm{def:g5:first-look-at-cells:cell}{cellules}, désigne un seul
individu (vrais jumeaux mis à part) dans toute l'humanité.
La même comparaison, menée au grossissement de la famille,
règle la filiation --- chaque position variable de l'enfant
doit pouvoir être distribuée à partir des jeux de ses deux
parents. La logique de l'album
(\cref{ch:g9:heredity-and-traits}) est devenue un
instrument.
\end{example}

\begin{example}[Transfusion et greffe]\label{ex:g9:unique-individuals:medicine}
La médecine négocie l'unicité tous les jours. Les groupes
sanguins (\cref{ex:g9:genes-and-dna:bloodgroups}) sont le
cas facile --- quatre grandes classes, appariables à partir
du registre. Les greffes d'organes rencontrent le cas
difficile~: le système immunitaire
(\cref{ch:g9:immune-defenses} en explique la machinerie) lit
des marqueurs du soi bien plus fins, assez uniques pour que
l'on cherche d'abord les donneurs dans la famille --- plus
la \omterm{prop:g6:classification-kinship:kinship}{parenté} est proche, plus la distribution est voisine.
\end{example}

\begin{proposition}[La variété est la réserve de l'espèce, mise en banque]\label{prop:g9:unique-individuals:reserve}
\Cref{rem:g8:sexual-reproduction:reserve} promettait le sens
du brassage~; les mécanismes l'encaissent aujourd'hui. Une
\omterm{def:g5:biodiversity:species}{espèce} faite d'individus uniques détient, éparpillés parmi
eux, des \omterm{def:g9:genes-and-dna:gene}{allèles} pour toutes les occasions --- la résistance
que certains portent à telle maladie, la tolérance que
d'autres portent à tel stress (le verger de
\cref{ex:g8:sexual-reproduction:bets}, à l'échelle de
l'\omterm{def:g5:biodiversity:species}{espèce}). Quand les \omterm{def:g6:exploring-our-environment:conditions}{conditions} tournent, les freins de
\cref{prop:g8:reproduction-and-environment:limits} tombent
inégalement sur cette variété --- et ce que fait cette
inégalité, génération après génération, est précisément le
sujet qu'ouvrent les deux prochains chapitres~: le
changement des \omterm{def:g5:biodiversity:species}{espèces}.
\end{proposition}
\admitted

\begin{method}[Passer au crible une affirmation «~génétique~»]\label{met:g9:unique-individuals:claims}
Pour tout titre de journal sur les \omterm{def:g9:genes-and-dna:gene}{gènes} et les gens~:
\begin{enumerate}
  \item vérification du niveau
        (\cref{met:g9:genes-and-dna:levels})~: l'affirmation
        porte-t-elle sur un \omterm{def:g9:genes-and-dna:gene}{gène}, sur un jeu, ou sur un
        motif familial~?
  \item vérification du tissage
        (\cref{prop:g9:heredity-and-traits:sorting})~:
        respecte-t-elle le tissage marge héréditaire et vie
        vécue, ou prétend-elle qu'un \omterm{def:g9:genes-and-dna:gene}{gène} égale un destin~?
  \item vérification de la variété~: traite-t-elle
        l'étalement d'une \omterm{def:g8:reproduction-and-environment:population}{population} pour la réserve qu'il
        est, ou comme des écarts à un jeu «~normal~» qui
        n'existe pas~?
  \item vérification des jumeaux~: survivrait-elle aux vrais
        jumeaux, qui partagent la distribution et diffèrent
        quand même~?
\end{enumerate}
\end{method}

\section{Exercices}

\begin{exercise}[$\star$]\label{exo:g9:unique-individuals:1}
Nomme les trois mécanismes du théorème de l'unicité.
\end{exercise}

\begin{exercise}[$\star$]\label{exo:g9:unique-individuals:2}
Qu'ajoute au brassage l'échange de segments pendant la
réduction de moitié~?
\end{exercise}

\begin{exercise}[$\star$]\label{exo:g9:unique-individuals:3}
Énonce les deux moitiés de
\cref{rem:g9:unique-individuals:alike} --- le pourcentage et
l'unicité --- en une seule phrase.
\end{exercise}

\begin{exercise}[$\star$]\label{exo:g9:unique-individuals:4}
Comment les vrais jumeaux finissent-ils malgré tout par
être distinguables~?
\end{exercise}

\begin{exercise}[$\star$]\label{exo:g9:unique-individuals:5}
À quelles deux questions la comparaison d'\omterm{def:g9:genes-and-dna:dna}{ADN} répond-elle
pour un tribunal~?
\end{exercise}

\begin{exercise}[$\star$]\label{exo:g9:unique-individuals:6}
Pourquoi cherche-t-on d'abord les donneurs de greffe dans
la famille~?
\end{exercise}

\begin{exercise}[$\star\star$]\label{exo:g9:unique-individuals:7}
Multiplie l'argument~: combine les espaces de demi-jeux des
deux parents et la loterie de la \omterm{def:g5:flower-to-fruit:steps}{fécondation} pour aboutir à
la conclusion «~jamais distribué auparavant~».
\end{exercise}

\begin{exercise}[$\star\star$]\label{exo:g9:unique-individuals:8}
Encaisse \cref{prop:g9:unique-individuals:reserve} sur le
vignoble de \cref{exo:g8:sexual-reproduction:10}~: que
manquait-il aux clones que détient une \omterm{def:g8:reproduction-and-environment:population}{population} née de
\omterm{prop:g3:flowers-fruits-seeds:transform}{graines}~?
\end{exercise}

\begin{exercise}[$\star\star$]\label{exo:g9:unique-individuals:9}
Applique \cref{met:g9:unique-individuals:claims} à~: «~un
\omterm{def:g9:genes-and-dna:gene}{gène} de la réussite scolaire découvert~».
\end{exercise}

\begin{exercise}[$\star\star$]\label{exo:g9:unique-individuals:10}
Pourquoi «~le jeu humain normal~» est-il une expression sans
référent~? Réponds par la réserve et par le théorème.
\end{exercise}

\begin{exercise}[$\star\star$]\label{exo:g9:unique-individuals:11}
Un test de filiation innocente un père mis en doute et
accusé à tort par le voisin de
\cref{exo:g9:heredity-and-traits:11}. Explique ce que
vérifie le test, position après position.
\end{exercise}

\begin{exercise}[$\star\star\star$]\label{exo:g9:unique-individuals:12}
«~Tous différents, tous parents, et les deux par le même
mécanisme.~» Écris le paragraphe de couronnement~: le
brassage, la \omterm{prop:g9:genes-and-dna:explains}{mutation} et le code partagé, chacun à sa place.
\end{exercise}

\section{Problème~: le séminaire sur une affaire non élucidée}

\begin{problem}[{Problème du week-end --- un cheveu, trois
questions, trente ans}]\label{pb:g9:unique-individuals:1}
Un séminaire (fictif) sur une affaire non élucidée~: un
cambriolage vieux de trente ans jamais résolu, un cheveu
conservé avec les \omterm{def:g5:first-look-at-cells:cell}{cellules} de sa \omterm{def:g1:plants-around-us:parts}{racine}, et les méthodes de
comparaison modernes. Le séminaire mène les élèves à travers
ce que l'identification par l'\omterm{def:g9:genes-and-dna:dna}{ADN} peut dire et ne peut pas
dire.

\medskip\noindent\textbf{Partie I --- Ce que contient le
cheveu.}
\begin{enumerate}
  \item Pourquoi la \emph{\omterm{def:g1:plants-around-us:parts}{racine}} du cheveu doit-elle être
        conservée pour une lecture complète
        (\cref{def:g6:cell-unit-of-life:parts} sait où
        habite l'\omterm{def:g9:genes-and-dna:dna}{ADN})~?
  \item Le laboratoire lit un ensemble de positions très
        variables. Pourquoi des positions variables ---
        qu'établiraient les positions partagées à 99 pour
        cent~?
  \item Le profil concorde avec un suspect à chaque
        position. Énonce ce qu'affirme cette concordance,
        et sa seule exception systématique
        (\cref{ex:g9:unique-individuals:twins}).
  \item La défense fait remarquer que le suspect a un vrai
        jumeau à l'étranger. Que doit faire le tribunal
        maintenant, et pourquoi~?
\end{enumerate}

\medskip\noindent\textbf{Partie II --- Les questions de
famille.}
\begin{enumerate}[resume]
  \item Un profil antérieur, partiel, avait désigné «~un
        proche parent de la famille K~». Explique comment
        une concordance \emph{partielle} se lit comme une
        \omterm{prop:g6:classification-kinship:kinship}{parenté} (la distribution de
        \cref{prop:g9:unique-individuals:theorem}, remontée
        à l'envers).
  \item Un juré demande si le profil révèle les \omterm{def:g9:heredity-and-traits:traits}{caractères}
        du suspect --- son allure, son tempérament, des
        «~\omterm{def:g9:genes-and-dna:gene}{gènes} du crime~». Applique
        \cref{met:g9:unique-individuals:claims} à la
        question, vérification par vérification.
  \item Un autre juré demande si la concordance ne pourrait
        pas être «~unique dans cette ville~» plutôt
        qu'unique dans l'humanité. Qu'est-ce qui fixe le
        nombre de positions lues~?
  \item L'animateur du séminaire avertit~: «~la \omterm{def:g1:living-or-not:living}{biologie}
        dit de qui sont les \omterm{def:g5:first-look-at-cells:cell}{cellules}~; l'affaire doit
        encore dire comment elles sont arrivées là.~»
        Distingue les deux affirmations.
\end{enumerate}

\medskip\noindent\textbf{Partie III --- La clôture du
séminaire.}
\begin{enumerate}[resume]
  \item Résume l'instrument en une phrase~: quels
        mécanismes, tirés de deux chapitres
        (\cref{ch:g9:chromosomes},
        \cref{ch:g9:genes-and-dna}), rendent un profil
        unique et une famille lisible.
  \item Pourquoi le même instrument marche-t-il sur un
        cheveu vieux de trente ans~? (Sur quelle propriété
        de l'\omterm{def:g9:genes-and-dna:dna}{ADN} --- la copie de
        \cref{prop:g9:genes-and-dna:explains} mise à part
        --- s'appuie-t-on~: la stabilité.)
  \item Le séminaire se termine sur la morale~: cite deux
        usages de l'unicité que ce chapitre montre au
        service des gens, et le mésusage historique contre
        lequel met en garde
        \cref{rem:g9:unique-individuals:alike}.
  \item Termine par le théorème réénoncé pour le prétoire~:
        qu'est-ce qui ne s'est jamais produit deux fois, et
        pourquoi.
\end{enumerate}
\end{problem}
